\documentclass[10pt]{extarticle}
\usepackage[a4paper,margin=1.5cm]{geometry}
\usepackage[T1]{fontenc}
\usepackage{amsmath,amssymb}
\usepackage{graphicx}
\usepackage[export]{adjustbox}
\usepackage{multicol}
\usepackage{enumitem}
\usepackage{booktabs,array,tabularx}
\setlength{\columnsep}{0.7cm}
\setlength{\parindent}{0pt}
\setlength{\parskip}{3pt}
\pagestyle{empty}
\makeatletter
% figure registry (figdefs.tex is generated by the extract step)
\newcommand{\qfigdef}[4]{\@namedef{qf@file@#1}{#2}\@namedef{qf@w@#1}{#3}\@namedef{qf@wide@#1}{#4}}
\newcommand{\qf@check}[1]{\@ifundefined{qf@file@#1}{\PackageError{score2}{Unknown figure id #1}{}}{}}
\newcommand{\qf@img}[4][]{\includegraphics[width=#2\dimexpr\@nameuse{qf@w@#4}pt\relax,max width=#3,max height=0.45\textheight,#1]{\@nameuse{qf@file@#4}}}
% Full-width questions: finalize wraps every question that contains a wide figure in \qfullbegin...\qfullend,
% so the WHOLE question (stem, figures, options) is set across the page in reading order, then the two columns
% resume. The question's head - its number, stem and first figure (\qfullsplit marks the end of the head) -
% is kept on one page, so a number is never left at a page foot with its figure overleaf; the rest may break
% across pages normally, so no big gaps are left. Without a \qfullsplit, the whole question is kept together
% if it is under 45% of a page. Never write these yourself.
\newif\ifqf@full
\newif\ifqf@boxed
\newsavebox{\qf@box}
\newcommand{\qf@open}{\global\qf@fulltrue\par\global\qf@boxedtrue
  \setbox\qf@box\vbox\bgroup\hsize\textwidth\linewidth\textwidth}
\newcommand{\qf@emit}[1]{\par\egroup\global\qf@boxedfalse
  \ifdim\dimexpr\ht\qf@box+\dp\qf@box\relax<#1\textheight\noindent\usebox{\qf@box}\par\else\unvbox\qf@box\fi}
\newcommand{\qfullsplit}{\ifqf@boxed\qf@emit{0.85}\fi}
\newcommand{\qf@close}{\ifqf@boxed\qf@emit{0.45}\fi\global\qf@fullfalse}
\newcommand{\qfullbegin}{\end{multicols}\qf@open}
\newcommand{\qfullend}{\qf@close\begin{multicols}{2}}
\newcommand{\qfullstart}{\qf@open}
\newcommand{\qfulldone}{\qf@close}
% \qfig[scale]{id}: block figure, centred, directly where it is written (size set by the tool).
% Fallback if a wide figure is not inside a full-width question: break out of the columns in place - no floats.
\newcommand{\qfig}[2][1]{\qf@check{#2}%
  \ifqf@full\qf@blk{#1}{#2}\else
  \ifnum\@nameuse{qf@wide@#2}=1 \expandafter\@firstoftwo\else\expandafter\@secondoftwo\fi
  {\end{multicols}\begin{center}\qf@img{#1}{\textwidth}{#2}\end{center}\begin{multicols}{2}}%
  {\qf@blk{#1}{#2}}\fi}
\newcommand{\qf@blk}[2]{\par\smallskip{\centering\qf@img{#1}{\linewidth}{#2}\par}\smallskip}
% \qfigi[scale]{id}: inline figure (inside option lists / tables / minipages); never breaks columns.
\newcommand{\qfigi}[2][1]{\qf@check{#2}\qf@img[valign=c]{#1}{\linewidth}{#2}}
\makeatother
% structure macros - lint reads these to check completeness and figure placement
\newcommand{\papersec}[1]{\par\bigskip{\centering\Large\bfseries PAPER-#1\par}\medskip}
\newcommand{\partsec}[2]{\par\medskip\noindent{\large\bfseries PART-#1: #2}\par\smallskip}
\newcommand{\sect}[1]{\par\smallskip\noindent\textbf{#1}\par}
\newcommand{\qn}[1]{\par\smallskip\noindent\textbf{#1.}~}
\newcommand{\paraq}[2]{\par\smallskip\noindent\textbf{Paragraph for Questions #1 and #2}\par}
\newenvironment{qoptions}{\begin{itemize}[leftmargin=1.8em,itemsep=1pt,topsep=2pt,label={}]}{\end{itemize}}
\newcommand{\opt}[2]{\item[(#1)] #2}
% figure registry (generated by score2_tools.py finalize)
\qfigdef{p1_phy_s1_q01_1}{TOAS - Test 10 (12 Jun 2022)-img/p1_phy_s1_q01_1.png}{129.6}{0}
\qfigdef{p1_phy_s1_q03_1}{TOAS - Test 10 (12 Jun 2022)-img/p1_phy_s1_q03_1.png}{84.6}{0}
\qfigdef{p1_phy_s1_q04_1}{TOAS - Test 10 (12 Jun 2022)-img/p1_phy_s1_q04_1.png}{153.9}{0}
\qfigdef{p1_phy_s1_q05_1}{TOAS - Test 10 (12 Jun 2022)-img/p1_phy_s1_q05_1.png}{108.7}{0}
\qfigdef{p1_phy_s1_q06_1}{TOAS - Test 10 (12 Jun 2022)-img/p1_phy_s1_q06_1.png}{109.7}{0}
\qfigdef{p1_phy_s1_q08_1}{TOAS - Test 10 (12 Jun 2022)-img/p1_phy_s1_q08_1.png}{93.1}{0}
\qfigdef{p1_phy_s2_para9_10_1}{TOAS - Test 10 (12 Jun 2022)-img/p1_phy_s2_para9_10_1.png}{93.3}{0}
\qfigdef{p1_phy_s3_q02_1}{TOAS - Test 10 (12 Jun 2022)-img/p1_phy_s3_q02_1.png}{77.3}{0}
\qfigdef{p1_phy_s3_q03_1}{TOAS - Test 10 (12 Jun 2022)-img/p1_phy_s3_q03_1.png}{98.5}{0}
\qfigdef{p1_phy_s3_q05_1}{TOAS - Test 10 (12 Jun 2022)-img/p1_phy_s3_q05_1.png}{62.0}{0}
\qfigdef{p1_phy_s3_q07_1}{TOAS - Test 10 (12 Jun 2022)-img/p1_phy_s3_q07_1.png}{96.2}{0}
\qfigdef{p1_phy_s3_q08_1}{TOAS - Test 10 (12 Jun 2022)-img/p1_phy_s3_q08_1.png}{131.6}{0}
\qfigdef{p1_chem_s1_q03_1}{TOAS - Test 10 (12 Jun 2022)-img/p1_chem_s1_q03_1.png}{170.0}{0}
\qfigdef{p1_chem_s1_q06_1}{TOAS - Test 10 (12 Jun 2022)-img/p1_chem_s1_q06_1.png}{81.3}{0}
\qfigdef{p1_chem_s1_q06_2}{TOAS - Test 10 (12 Jun 2022)-img/p1_chem_s1_q06_2.png}{83.2}{0}
\qfigdef{p1_chem_s1_q06_4}{TOAS - Test 10 (12 Jun 2022)-img/p1_chem_s1_q06_4.png}{24.5}{0}
\qfigdef{p1_chem_s1_q06_3}{TOAS - Test 10 (12 Jun 2022)-img/p1_chem_s1_q06_3.png}{88.6}{0}
\qfigdef{p1_chem_s1_q07_1}{TOAS - Test 10 (12 Jun 2022)-img/p1_chem_s1_q07_1.png}{221.8}{0}
\qfigdef{p1_chem_s1_q08_1}{TOAS - Test 10 (12 Jun 2022)-img/p1_chem_s1_q08_1.png}{245.2}{0}
\qfigdef{p1_chem_s1_q08_2}{TOAS - Test 10 (12 Jun 2022)-img/p1_chem_s1_q08_2.png}{25.2}{0}
\qfigdef{p1_chem_s1_q08_3}{TOAS - Test 10 (12 Jun 2022)-img/p1_chem_s1_q08_3.png}{22.5}{0}
\qfigdef{p1_chem_s2_para9_10_1}{TOAS - Test 10 (12 Jun 2022)-img/p1_chem_s2_para9_10_1.png}{193.7}{0}
\qfigdef{p1_chem_s3_q07_1}{TOAS - Test 10 (12 Jun 2022)-img/p1_chem_s3_q07_1.png}{163.1}{0}
\qfigdef{p1_chem_s3_q07_2}{TOAS - Test 10 (12 Jun 2022)-img/p1_chem_s3_q07_2.png}{157.3}{0}
\qfigdef{p1_chem_s3_q07_3}{TOAS - Test 10 (12 Jun 2022)-img/p1_chem_s3_q07_3.png}{164.1}{0}
\qfigdef{p1_chem_s3_q07_4}{TOAS - Test 10 (12 Jun 2022)-img/p1_chem_s3_q07_4.png}{154.2}{0}
\raggedcolumns
\makeatletter\@beginparpenalty=10000\makeatother
\newcolumntype{L}[1]{>{\raggedright\arraybackslash}p{#1}}
\begin{document}
\begin{center}{\Large\bfseries TOAS --- Test 10 (12 Jun 2022)}\end{center}
\begin{multicols}{2}
\raggedcolumns
\partsec{1}{PHYSICS}
\sect{SECTION-I(i)}
\qn{1} Figure shows a potentiometer circuit. The length of potentiometer wire AB is 120 cm and its resistance is 2$\Omega$. The internal resistance of 2V cell is 1$\Omega$. Initially the switches S$_1$ and S$_2$ are open and S$_3$ is closed.
\qfig{p1_phy_s1_q01_1}
\begin{qoptions}
\opt{A}{When only S$_2$ is open, balanced length $\ell=20$ cm.}
\opt{B}{When only S$_1$ is open, balanced length $\ell=60$ cm}
\opt{C}{When all three switches are closed, balanced length is 60 cm.}
\opt{D}{When S$_3$ is open, S$_2$ is also open but S$_1$ is closed, balance length is 60 cm.}
\end{qoptions}
\qn{2} A thin rod of length $\ell$ is kept vertically (along y-axis) with end A being at rest on a smooth horizontal plane. A slight disturbance makes the end A slip on the plane (along x-axis) and the rod falls. Then,
\begin{qoptions}
\opt{A}{the path described by a point P at a distance of $\dfrac{\ell}{4}$ from A is a quadrant of a circle.}
\opt{B}{the centre of mass of the rod moves along a straight line, parallel to y-axis.}
\opt{C}{the acceleration of the centre of mass as the rod becomes horizontal is $\dfrac{3g}{4}$.}
\opt{D}{the greatest angular velocity of the rod about the instantaneous centre of rotation is $\sqrt{\dfrac{3g}{\ell}}$.}
\end{qoptions}
\qn{3} A, B and C are the vertices of a uniform equilateral triangular thin lamina of mass M. Three particle of masses $\dfrac M3$, each are located at the midpoints P, Q and R of the edges of the lamina. For this system, (each edge of lamina = a)
\qfig{p1_phy_s1_q03_1}
\begin{qoptions}
\opt{A}{the moment of inertia of the system about an axis through the centroid O perpendicular to the plane is $\dfrac{Ma^2}{6}$.}
\opt{B}{the moment of inertia of the system about an axis through O perpendicular to the plane is $\dfrac{5Ma^2}{12}$}
\opt{C}{the moment of inertia of the system about an axis through A perpendicular to the plane is $\dfrac{13Ma^2}{12}$}
\opt{D}{the moment of inertia of the system about an axis through C perpendicular to the plane is $\dfrac{5Ma^2}{6}$}
\end{qoptions}
\qn{4} A network of four unequal resistors $R_1$, $R_2$, $R_3$ and $R_4$ are connected in a network (see diagram) that contains a source of e.m.f. E and an ammeter A. The source has negligible internal resistance. For this network.
\qfig{p1_phy_s1_q04_1}
\begin{qoptions}
\opt{A}{If $R_3>R_4$, the current in the resistor $R_3$ will be less than that shown by the ammeter.}
\opt{B}{If $R_3<R_4$, the current in the resistor $R_2$ will be greater than that shown by the ammeter.}
\opt{C}{If the source of e.m.f. and the ammeter were physically interchanged in the network the ammeter will show the same current.}
\opt{D}{If the resistors $R_2$ and $R_1$ were physically interchanged, the ammeter will show the same current}
\end{qoptions}
\qn{5} In the given network, the powers generated in the resistors $R_1$ and $R_2$ are $P_1$ and $P_2$ and the power generated in R is P. For this network,
\qfig{p1_phy_s1_q05_1}
\begin{qoptions}
\opt{A}{If $P_1=P_2$, it necessarily follows that $R_1=R_2$}
\opt{B}{If $R=R_1+R_2$, then $P>P_1+P_2$}
\opt{C}{If $P=P_1+P_2$, then $R\leqslant\dfrac{R_1+R_2}{4}$}
\opt{D}{If $R_1=R_2$, then it necessarily follows that $P=P_1+P_2$}
\end{qoptions}
\qn{6} A charge q is placed at a distance r from centre of conducting sphere as shown in figure. Now sphere is earthed then correct options are :-
\qfig{p1_phy_s1_q06_1}
\begin{qoptions}
\opt{A}{Induced charge on sphere is $-\left(\dfrac Rr\right)q$}
\opt{B}{Induced charge on sphere is $-q$}
\opt{C}{Force acting between sphere and external charge q is $\dfrac{1}{4\pi\varepsilon_0}\dfrac{q^2Rr}{(r^2-R^2)^2}$}
\opt{D}{Force between sphere and external charge q is $\dfrac{q^2}{r^2}$.}
\end{qoptions}
\qn{7} A uniform disc of mass 2kg and radius 1m is mounted on a horizontal axle supported on fixed frictionless bearings. A light chord is wrapped around the rim of the disc and mass of 1kg is tied to the free end. If it is released from rest (take $g=10$ m/s$^2$)
\begin{qoptions}
\opt{A}{the tension in the chord is 5N.}
\opt{B}{in first four seconds the angular displacement of the disc is 40 rad.}
\opt{C}{the work done by the torque on the disc in first four seconds is 200 J.}
\opt{D}{the increase in the kinetic energy of the mass in the first four seconds is 200 J.}
\end{qoptions}
\qn{8} A uniform rod of length $\ell$ is falling down with a velocity $v_0$ when one of it's end hits a fixed edge as shown. (Consider impulsive force to be acting perpendicular to the length of rod)
\qfig{p1_phy_s1_q08_1}
\begin{qoptions}
\opt{A}{If the collision is elastic, the centre of mass may have upwards velocity just after the collision.}
\opt{B}{If the collision is inelastic, the centre of mass may come to rest just after the collision.}
\opt{C}{If the collision is perfectly inelastic, the centre of mass has a downward velocity just after the collision.}
\opt{D}{If the collision is elastic, the centre of mass has a downward velocity just after the collision.}
\end{qoptions}
\sect{SECTION-I(ii)}
\paraq{9}{10}
A small particle of mass $m$ is fixed to the perimeter of a ring of same mass and radius `r' and ring is placed on a horizontal plane. Friction is negligible on horizontal plane. Initially particle is at top most point, then the system is released from rest. Answer next two questions when the particle is at the same height as the centre of ring after being released from topmost point. Assume that the ring stays in vertical plane during its motion under consideration.
\qfig{p1_phy_s2_para9_10_1}
\qn{9} Mark the correct option(s) :-
\begin{qoptions}
\opt{A}{Speed of centre of mass of system is $\sqrt{\dfrac{gr}{4}}$.}
\opt{B}{Angular velocity of system is $\omega=\sqrt{\dfrac gr}$}
\opt{C}{Contact point of ring with ground is instantaneously at rest.}
\opt{D}{Centre of ring is instantaneously at rest.}
\end{qoptions}
\qn{10} Mark the correct option(s) :
\begin{qoptions}
\opt{A}{Instantaneous acceleration of the centre of ring is $\dfrac g2$.}
\opt{B}{Acceleration of centre of ring initially points towards left and then towards right.}
\opt{C}{The force exerted by the ring on the ground at the instant given is $\dfrac32mg$.}
\opt{D}{The force exerted by the ring on the ground at the instant given is $2mg$.}
\end{qoptions}
\paraq{11}{12}
A rocket is launched from and returns to a spherical planet of radius $R$ in such a way that its velocity vector on return is parallel to its launch vector. The angular separation at the centre of the planet between the launch and arrival points is $\theta$.
\qn{11} Which of the following can be the distance of the rocket above the surface of the planet ?
\begin{qoptions}
\opt{A}{$R$}
\opt{B}{$\dfrac R2$}
\opt{C}{$\dfrac R3$}
\opt{D}{$4R$}
\end{qoptions}
\qn{12} How long does the flight of rocket take, if the period of satellite flying around the planet just above its surface is $T_0$ ?
\begin{qoptions}
\opt{A}{$T_0\left(\dfrac12+\dfrac1\pi\right)$ if $\theta\to0$}
\opt{B}{$T_0\left(\dfrac12+\dfrac{1}{\sqrt2\pi}\right)$ if $\theta=90^\circ$}
\opt{C}{$T_0\left(\dfrac12+\dfrac2\pi\right)$ if $\theta\to0$}
\opt{D}{$T_0\left(\dfrac12+\dfrac{\sqrt2}{\pi}\right)$ if $\theta=90^\circ$}
\end{qoptions}
\sect{SECTION-II}
\qn{1} A particle is uncharged and is thrown vertically upward from ground level with a speed of $5\sqrt5$ m/s. As a result, it attains a maximum height $h$. The particle is then given a positive charge $+q$ and reaches the same maximum height $h$ when thrown vertically upward with a speed of 13 m/s. Finally, the particle is given a negative charge $-q$. Ignoring air resistance, determine the speed (in m/s) with which the negatively charged particle must be thrown vertically upward, so that it attains exactly the same maximum height $h$.
\qn{2} A smooth parabolic wire track defined by relation $y=\dfrac{x^2}{2a}$ is fixed in x-y plane in a gravity free space as shown in figure. An electric field $\vec E=\alpha y\hat i+\alpha x\hat j$ (where $\alpha$ is a constant) exist in the region. A small ring the mass `m', charge `q', which just fits in wire track, is released from rest at position $A\left(-a,\dfrac a2\right)$. Find speed of the ring when it reaches at point $B\left(a,\dfrac a2\right)$ if $q\alpha=\dfrac{2m}{a^2}$.
\qfig{p1_phy_s3_q02_1}
\qn{3} Consider four concentric conducting spherical shells of radius $a$, $2a$, $3a$ and $4a$ having charges $Q$, $-Q$, $2Q$ and $Q$ respectively as shown in figure. Total electrostatic potential energy of this charge configuration is $\left(\dfrac{nQ^2}{\pi\varepsilon_0a}\right)$. Find $n$.
\qfig{p1_phy_s3_q03_1}
\qn{4} A voltmeter coil has resistance 50.0 $\Omega$ and a resistance of 1.15 k$\Omega$ is connected in series. It can read potential differences up to 12 volts. If this same coil is used to construct an ammeter which can measure currents upto 2.0 A, what should be the resistance of the shunt used ?
\qn{5} A particle of mass $m$ is fired from the point A (as shown in the figure) situated at a distance of $2R$ from the centre of the planet with a velocity $\dfrac{1}{\sqrt2}$ times the escape velocity from point A. The radius of the planet is $R$. The maximum distance of the particle from the centre in the subsequent motion is $nR$. Find the value of $n$.
\qfig{p1_phy_s3_q05_1}
\qn{6} A metal sphere of radius $a$ is surrounded by a concentric metal spherical shell of radius $b$($b>a$). The space between the spheres is filled with material whose electrical conductivity $\sigma$ varies with the electric field strength $E$ according to the relation $\sigma=KE$ where $K$ is constant. A potential difference $V$ is maintained between inner sphere and outer shell. If the current (in ampere) between the inner sphere and outer shell is $x$, then find $x/4$. [Here $K=\dfrac1\pi\times10^{-2}\,\mathrm{ohm^{-1}volt^{-1}}$, $V=20$ volts, $b=ea$ ($e$ = exponential) and $\pi=22/7$]
\qn{7} In the figure ammeter (I) reads a current of 10 mA, while the voltmeter reads a potential difference of 3V. What does ammeter (II) (in mA) read ? The ammeters are identical, the internal resistance of the battery is negligible. (Consider all ammeters and voltmeters as non ideal]
\qfig{p1_phy_s3_q07_1}
\qn{8} A disc of radius `5cm' rolls on a horizontal surface with linear velocity $v=1\hat i$ m/s and angular velocity 50 rad/sec. Height of particle from ground on rim of disc which has velocity in vertical direction is (in cm) :
\qfig{p1_phy_s3_q08_1}
\raggedcolumns
\partsec{2}{CHEMISTRY}
\sect{SECTION-I(i)}
\qn{1} For non-stoichiometric cuprous oxide, Cu$_{1.80}$ O$_{1.00}$, select the correct information(s). (Cu = 63.5) :
\begin{qoptions}
\opt{A}{It behaves as p-type semiconductor}
\opt{B}{\% of total copper as Cu(I) is 11.11\%}
\opt{C}{It has metal deficiency defect}
\opt{D}{Density of solid is 20\% less than the ideal Cu$_2$O crystal}
\end{qoptions}
\qn{2} Liquids `A' and `B' form an ideal solution. At 25$^\circ$C, $P_A^\circ=25$ cm Hg and $P_B^\circ=75$ cm Hg. If 5.0 moles of liquid `A' is mixed with 20.0 moles of liquid `B' and the pressure over the system is maintained at 50 cm Hg, the correct information(s) regarding the final composition of system is/are :
\begin{qoptions}
\opt{A}{Vapour contains 5.0 moles of `A'.}
\opt{B}{Liquid solution contains 50\% of `A', by moles.}
\opt{C}{Vapours contains 75\% of `B', by moles}
\opt{D}{50\% of the total moles of liquid taken will vaporise.}
\end{qoptions}
\qn{3} \qfig{p1_chem_s1_q03_1}
If above setup is used during electro-magnetic separation of magnetic particles with non-magnetic particles then which of the following observation will be observed.
\begin{qoptions}
\opt{A}{Most of the part of magnetic ore will be collected in container `B'}
\opt{B}{Silicates impurities are collected in container `A'}
\opt{C}{It is based on physical properties of matter}
\opt{D}{In case of casseterite, gangue particles are collected in container `A'}
\end{qoptions}
\qn{4} Which of the following is/are correct for commercial extraction of metals?
\begin{qoptions}
\opt{A}{During Matte formation FeO is removed in the form of slag}
\opt{B}{Copper from it's low grade ore can be extracted by hydrometallurgical reduction using H$_2$ gas.}
\opt{C}{Almost all Au is extracted by amalgamation process completely from Aluvial sand}
\opt{D}{Metal sulphides are more stable as compared to corresponding metal oxides}
\end{qoptions}
\qn{5} Three complex compounds I, II and III are $[\mathrm{M(H_2O)_6}]^{3+}$, $[\mathrm{M(en)_3}]^{3+}$, $[\mathrm{M(CN)_6}]^{3-}$ respectively. Complex I and III have green and red colour respectively then identify the correct possible matching colour for complex II.
\begin{qoptions}
\opt{A}{Blue}
\opt{B}{Orange}
\opt{C}{Yellow}
\opt{D}{Violet}
\end{qoptions}
\qn{6} (P) + NaHSO$_3\rightleftharpoons$ (Q)\par
Compound (P) having molecular formula `C$_3$H$_6$O' \& gives (+ve) haloform test. Choose the correct statement(s)?
\begin{qoptions}
\opt{A}{(Q) is \qfigi{p1_chem_s1_q06_1} and crystalline in nature}
\opt{B}{(Q) is \qfigi{p1_chem_s1_q06_2} and non-crystalline in nature}
\opt{C}{(P) can be aldehyde, methyl ketone \& some cyclic ketones as \qfigi{p1_chem_s1_q06_4}}
\opt{D}{(Q) is \qfigi{p1_chem_s1_q06_3} and crystalline in nature}
\end{qoptions}
\qn{7} \qfig{p1_chem_s1_q07_1}
Correct statement(s) is/are
\begin{qoptions}
\opt{A}{(A) is vicinol diol}
\opt{B}{(B) is a dicarboxylic acid}
\opt{C}{(B) gives Tollen's test}
\opt{D}{(D) gives Iodoform test}
\end{qoptions}
\qn{8} \qfig{p1_chem_s1_q08_1}
Correct statement is/are :
\begin{qoptions}
\opt{A}{Compound F is \qfigi{p1_chem_s1_q08_2}}
\opt{B}{Compound E is \qfigi{p1_chem_s1_q08_3}}
\opt{C}{A is acetophenone}
\opt{D}{In reaction from A to B, 3 equivalent of I$_2$ consumed, with respect to 1 equivalent of A.}
\end{qoptions}
\sect{SECTION-I(ii)}
\paraq{9}{10}
A substance `X' converts into `Z' through some intermediate, `Y'.
\qfig{p1_chem_s2_para9_10_1}
\qn{9} The correct concentration(s) of `X', `Y' and `Z' at $t=20$ min is ($e^{-1}=0.36$, $e^{-10}=4.5\times10^{-5}$) :
\begin{qoptions}
\opt{A}{[X] = 0.18 M}
\opt{B}{[Y] $\approx$ 0.02 M}
\opt{C}{[Z] $\approx$ 0.30 M}
\opt{D}{[X] + [Y] + [Z] = 0.50 M}
\end{qoptions}
\qn{10} The correct information(s) is/are ($\ell$n 2 = 0.7, $\ell$n = 1.6) :
\begin{qoptions}
\opt{A}{The concentration of `Y' becomes maximum at about 5.11 min. from the start of reaction}
\opt{B}{The rate of formation of `Z' becomes maximum at about 5.11 min. from the start of reaction}
\opt{C}{The rate of disappearance of `X' becomes equal to the rate of appearance of `Z' at about 5.11 min. from the start of reaction}
\opt{D}{Concentration of `X' becomes 0.25 M at about 5.11 min. from the start of reaction}
\end{qoptions}
\paraq{11}{12}
An octahedral homoleptic high spin complex of `3d' series central metal ion having 3d$^n$ electronic configuration have $t_{2g}^{x}e_g^2$ arrangement of electrons.\par
Answer the following based on above paragraph.
\qn{11} Choose the correct pair of central atom/ion and ligand for the above mentioned complex.
\begin{qoptions}
\opt{A}{Fe$^{2+}$, NH$_3$}
\opt{B}{Co$^{3+}$, NH$_3$}
\opt{C}{Fe$^{3+}$, F$^{-}$}
\opt{D}{Ni$^{2+}$, H$_2$O}
\end{qoptions}
\qn{12} For the maximum possible value of `x' the value of spin only magnetic moment will be.
\begin{qoptions}
\opt{A}{$\sqrt8$ B.M.}
\opt{B}{$\sqrt{15}$ B.M.}
\opt{C}{$\sqrt{24}$ B.M.}
\opt{D}{$\sqrt{35}$ B.M.}
\end{qoptions}
\sect{SECTION-II}
\qn{1} N$_2$ and H$_2$ gases are mixed in 1 : 3 mole ratio. At equilibrium at 100 atm, the equilibrium mixture is found to contain 20 mole percent NH$_3$. For the reaction :\par
N$_2$(g) + 3H$_2$(g) $\rightleftharpoons$ 2NH$_3$(g)\par
$K_p$ is $\left(\dfrac{x}{216}\right)$ atm$^{-2}$. The value of `x' is.
\qn{2} 2 moles of an ideal gas undergoes reversible isothermal expansion from 5 dm$^3$ is 20 dm$^3$ at 500 K. The amount of heat absorbed by the gas (in Kcal) is ($\ell$n 2 = 0.693).
\qn{3} The standard enthalpies of combustion of C(graphite), H$_2$(g) and C$_3$H$_8$(g) are $-393$, $-285$ and $-3280$ kJ/mole, respectively. The standard enthalpy of formation (in KJ/mole) of C$_3$H$_8$(g) is.
\qn{4} If all the bond angles w.r.t. `N' and `B' are assumed as 120$^\circ$ in BH$_2$(NMe$_2$) molecule then maximum number of atoms in one plane will be :-
\qn{5} Some of the following metallurgical process are given below :-\par
a. Electro magnetic separation\par
b. Froth floatation process\par
c. Calcination\par
d. Roasting\par
e. Leaching\par
f. Reduction by coke\par
g. Self reduction\par
h. Electrolytic refining\par
i. refining of metal by liquation\par
The number of above stages involved in the extraction of Zn from Calamine and Cu from Chalcosite are X and Y. the value of X + Y is.
\qn{6} R-2-chlorohexane $\xrightarrow[h\nu]{\mathrm{Cl_2}}$ Number of monochloro products.
\qn{7} Find number of reaction correctly matched, with their major product (or statement written with them) :
\qfig{p1_chem_s3_q07_1}
(4) CH$_3$CHO gives positive Tollen's test as well as Lucas test
\qfig{p1_chem_s3_q07_2}
\qfig{p1_chem_s3_q07_3}
\qfig{p1_chem_s3_q07_4}
\qn{8} Which of the following compound when reacts with CH$_3$MgCl gives methane gas\par
(i) Tertiary butylchloride\ \ (ii) Tertiary butylalcohol\ \ (iii) Di-butylether\ \ (iv) Phenol\par
(v) Pyrole\ \ (vi) Pyridine
\raggedcolumns
\partsec{3}{MATHEMATICS}
\sect{SECTION-I(i)}
\qn{1} $\displaystyle\lim_{x\to\infty} x^2\tan\left(\ln\sqrt{\cos\dfrac{\pi}{x}}\right)$ is equal to
\begin{qoptions}
\opt{A}{$-\dfrac{\pi^2}{2}$}
\opt{B}{$\dfrac{\pi^2}{2}$}
\opt{C}{$-\dfrac{\pi^2}{4}$}
\opt{D}{$\dfrac{\pi^2}{4}$}
\end{qoptions}
\qn{2} If $\displaystyle\int\dfrac{x^3-8}{x^3\sqrt{x^4-x^3+2}}\,dx=\dfrac{k\sqrt{x^4-x^3+2}}{x^2}+C$, then the value of $k$ is
\begin{qoptions}
\opt{A}{3}
\opt{B}{2}
\opt{C}{7}
\opt{D}{4}
\end{qoptions}
\qn{3} If $\displaystyle\int\dfrac{3z^3-8z+5}{\sqrt{z^2-4z-7}}\,dz=(z^2+az+36)\sqrt{z^2-4z-7}+b\ln\left|z-2+\sqrt{z^2-4z-7}\right|+C$, where $a,b\in I$, then
\begin{qoptions}
\opt{A}{$a>b$}
\opt{B}{$a<b$}
\opt{C}{$a+b=117$}
\opt{D}{exactly one of $a,b$ is prime}
\end{qoptions}
\qn{4} Let $f(x)=\cos^{-1}\left(\ln\dfrac{x+1}{x}\right)$, then
\begin{qoptions}
\opt{A}{domain of $f(x)$ contains $(-\infty,-1)\cup(1,\infty)$}
\opt{B}{range of $f(x)$ contains $\left[0,\dfrac{1}{e-1}\right]$}
\opt{C}{domain of $f(x)$ is $\left(-\infty,\dfrac{e}{1-e}\right]\cup\left[\dfrac{1}{e-1},\infty\right)$}
\opt{D}{range of $f(x)$ is $(-\infty,-1]\cup(1,\infty)$}
\end{qoptions}
\qn{5} Let $f(x)$ be a twice differentiable function such that $f(2+x)=f(2-x)$ and $f'\left(\dfrac12\right)=0=f'(1)$. Then the minimum number of roots of $f''(x)=0$ in $(0,4)$ is
\begin{qoptions}
\opt{A}{4}
\opt{B}{3}
\opt{C}{2}
\opt{D}{5}
\end{qoptions}
\qn{6} Let $F(n,x)=\displaystyle\sum_{r=1}^{n}\tan^{-1}\left(\dfrac{2x}{4+r(r+1)x^2}\right)$ and $\dfrac{d}{dx}F(n,x)=G(n,x)$, then
\begin{qoptions}
\opt{A}{$F(8,2)=\tan^{-1}\dfrac45$}
\opt{B}{$\displaystyle\lim_{n\to\infty}F(n,2\sqrt3)=\dfrac\pi6$}
\opt{C}{if $G\left(n,\dfrac13\right)<0$ then $n\ge36$}
\opt{D}{$\displaystyle\lim_{n\to\infty}G(n,1)=-\dfrac25$}
\end{qoptions}
\qn{7} If $a+b+c=4$ and $a^2+b^2+c^2+3(ab+bc+ac)=21$, where $a,b,c\in\mathbb{R}$, then
\begin{qoptions}
\opt{A}{minimum value of $(a+b)(b+c)(a+c)$ is 18}
\opt{B}{minimum value of $(a+b)(b+c)(a+c)$ is 24}
\opt{C}{maximum value of $(a+b)(b+c)(a+c)$ is $\dfrac{490}{27}$}
\opt{D}{maximum value of $(a+b)(b+c)(a+c)$ is $\dfrac{471}{27}$}
\end{qoptions}
\qn{8} Let $f(x)=\begin{cases}x[x], & -4\le x<-1\\ 1-x, & -1\le x<1\\ x^2-3x+2, & 1\le x<2\\ (x-2)[x], & 2\le x\le4\end{cases}$ where $[\cdot]$ denotes the greatest integer function. Then
\begin{qoptions}
\opt{A}{number of points of non-differentiability of $f(x)$ is 6}
\opt{B}{sum of global maximum and global minimum values of $f(x)$ is $\dfrac{63}{4}$}
\opt{C}{$f(x)$ has 3 points of extrema}
\opt{D}{number of solutions of $f(x)=4$ is 3}
\end{qoptions}
\sect{SECTION-I(ii)}
\paraq{9}{10}
Let $f(x)$ be a differentiable function such that $f(x)=\dfrac12\left[f(xy)+f\left(\dfrac xy\right)\right]$ for all $x,y\in\mathbb{R}^+$, $f(1)=0$ and $f'(1)=2$.
\qn{9} $f(x)-f(y)$ is equal to
\begin{qoptions}
\opt{A}{$f\left(\dfrac yx\right)$}
\opt{B}{$f\left(\dfrac xy\right)$}
\opt{C}{$f(2x)$}
\opt{D}{$f(2y)$}
\end{qoptions}
\qn{10} $f'(3)$ is
\begin{qoptions}
\opt{A}{$\ge\dfrac13$}
\opt{B}{$\ge\dfrac23$}
\opt{C}{$\ge\dfrac12$}
\opt{D}{$\ge\dfrac14$}
\end{qoptions}
\paraq{11}{12}
Let $f(x)=ax^3+x^2-(2a+1)x+2$, $a\ge0$.
\qn{11} The number of values of $a$ for which Rolle's theorem is applicable in $[0,1]$ is $k$, then $k$ can be
\begin{qoptions}
\opt{A}{$\le0$}
\opt{B}{1}
\opt{C}{2}
\opt{D}{4}
\end{qoptions}
\qn{12} If $k_1$ and $k_2$ are the local maximum and local minimum values of $f(x)$ respectively, then the sum of roots of $f(x)=k_1$ and $f(x)=k_2$ at which $f'(x)\ne0$ is
\begin{qoptions}
\opt{A}{$-\dfrac{2}{3a}$}
\opt{B}{$-\dfrac1a$}
\opt{C}{$-\dfrac2a$}
\opt{D}{$-\dfrac{1}{3a}$}
\end{qoptions}
\sect{SECTION-II}
\qn{1} If $\displaystyle\int\dfrac{x^{17}}{(1+x^2)^{10}}\,dx=\dfrac1n\left(\dfrac{x^2}{1+x^2}\right)^m+C$, where $m,n\in N$, then find $\dfrac mn$ in lowest form.
\qn{2} A circle of radius $2\sqrt{10}$ cm is divided into two segments by a chord of length 12 cm. If the maximum area of a rectangle inscribed in the smaller segment is $\lambda$, then find $\dfrac{\lambda}{\sqrt{15}}$.
\qn{3} Let $f(x)=\dfrac{195x(x^2+4)}{9x^4+97x^2+144}$. If the range of $y=|f(x)|$ is $\left[0,\dfrac ab\right]$, where $a$ and $b$ are relatively prime, then find the value of $a-4b$.
\qn{4} If $\displaystyle\int\dfrac{dx}{(\sec x+\operatorname{cosec}x+\tan x+\cot x)^2}=\dfrac x\alpha+\dfrac{\cos\left(x+\frac\pi4\right)}{\beta}+\dfrac{\cos2x}{\gamma}+\delta$, then find $|\alpha+\beta+\gamma|$.
\qn{5} Let $f(x)=\max\{\min(x,x^2),\max(x,x^2)\}$, $x\ge0$. Find the value of $\dfrac{d}{dx}\left((f\circ f)^{-1}(x)\right)$ at $x=16$.
\qn{6} Find the value of $\displaystyle\lim_{\theta\to0^+}(\sin\theta)^{\sin\theta-\sin^2\theta}$.
\qn{7} If $\dfrac{d^3y}{dx^3}=-\left(\dfrac{dy}{dx}\right)^2\left[k\dfrac{d^2x}{dy^2}\cdot\dfrac{d^2y}{dx^2}+\left(\dfrac{dy}{dx}\right)^2\dfrac{d^3x}{dy^3}\right]$, then find the value of $k$.
\qn{8} Let $d=(x_2-x_1)^2+\left(-x_2-4+2\sqrt{x_1+2}\right)^2$, where $x_1>-2$ and $x_2\in\mathbb{R}$. Find the minimum value of $d$.
\end{multicols}
\end{document}
